\section{Related Work}

\subsection{Mixture of Experts}

\textbf{Switch Transformer.} Fedus et al. (2022) simplify MoE to top-1 routing, 
reducing communication overhead while maintaining capacity. BAM-UMoE uses 
top-2 routing, which has been shown to give better quality per FLOP.

\textbf{GShard.} Lepikhin et al. (2021) introduce top-2 gating with auxiliary 
load-balancing loss. BAM-UMoE uses the same formulation for its MoE branch.

\textbf{Expert Choice.} Zhou et al. (2022) reverse the routing direction: experts 
choose tokens instead of tokens choosing experts. Orthogonal to our approach.

\textbf{Soft MoE.} Puigcerver et al. (2023) use soft routing without discrete 
top-k. BAM-UMoE uses hard top-2 for simplicity.

\subsection{Post-Hoc MoE}

\textbf{MoEfication.} Zhang et al. (2022) split pretrained FFN layers into 
multiple experts without retraining. BAM-UMoE is orthogonal -- it adds new 
capacity rather than partitioning existing capacity.

\textbf{Sparse Upcycling.} Komatsuzaki et al. (2023) initialize MoE from a dense 
model by duplicating FFN weights. BAM-UMoE does not duplicate; it adds 
a zero-initialized MoE branch.

\textbf{Lifelong MoE.} Chen et al. (2023) adapt MoE to new tasks while preserving 
old ones. Different motivation.

\subsection{Parameter-Efficient Fine-Tuning}

\textbf{LoRA.} Hu et al. (2021) parameterize a weight update with a low-rank 
product $BA$, $B = 0$ at initialization. BAM-UMoE uses LoRA for base 
attention adaptation alongside the MoE branch.

\textbf{AdaLoRA.} Zhang et al. (2023) allocate LoRA rank adaptively. Orthogonal.

\textbf{PiSSA.} Meng et al. (2024) initialize LoRA from principal components. 
Compatible with BAM-UMoE.

\subsection{Attention Augmentation}

\textbf{Attention Refinement.} Many works augment attention with additional 
branches. BAM-UMoE's specific contribution is applying \textbf{MoE} to the 
value pathway.

\textbf{Gated Attention.} Gating mechanisms in attention have been studied 
(Parisotto et al., 2020; Bapna et al., 2021). BAM-UMoE uses gating to 
control the MoE contribution's ramp-up.

\subsection{Position of Our Work}

\begin{table}[H]
\centering
\small
\begin{tabular}{lll}
\toprule
\textbf{Feature} & \textbf{Prior MoE} & \textbf{BAM-UMoE} \\
\midrule
Application & Pretraining & Post-hoc \\
Step-0 preservation & Not typically & \textbf{Explicit} \\
Base weights & Trained & Frozen \\
Target & FFN capacity & \textbf{Value pathway} \\
LoRA integration & Rare & Yes \\
\bottomrule
\end{tabular}
\end{table}

BAM-UMoE is unique in explicitly targeting the \textbf{value pathway} with MoE 
and providing a \textbf{three-part step-0 preservation} mechanism.
\begin{lstlisting}

---
\end{lstlisting}